\begin{ex}
Số liền sau của số 999 là số nào?
\shortans{1000}
\loigiai{
Số liền sau của 999 là $999 + 1 = 1000$.
}
\end{ex}

\begin{ex}
Số liền trước của số 500 là số nào?
\shortans{499}
\loigiai{
Số liền trước của 500 là $500 - 1 = 499$.
}
\end{ex}

\begin{ex}
Số gồm 6 trăm, 3 chục và 8 đơn vị viết là gì?
\shortans{638}
\loigiai{
Số gồm 6 trăm, 3 chục, 8 đơn vị là 638.
}
\end{ex}

\begin{ex}
Số 705 có chữ số hàng chục là bao nhiêu?
\shortans{0}
\loigiai{
Trong số 705, chữ số hàng chục là 0.
}
\end{ex}

\begin{ex}
Số lớn nhất có ba chữ số khác nhau là số nào?
\shortans{987}
\loigiai{
Số lớn nhất có 3 chữ số khác nhau là 987.
}
\end{ex}

\begin{ex}
Số bé nhất có ba chữ số khác nhau là số nào?
\shortans{102}
\loigiai{
Số bé nhất có 3 chữ số khác nhau là 102.
}
\end{ex}

\begin{ex}
Số tròn trăm liền sau của số 400 là số nào?
\shortans{500}
\loigiai{
Số tròn trăm liền sau của 400 là 500.
}
\end{ex}

\begin{ex}
Số tròn chục liền trước của số 320 là số nào?
\shortans{310}
\loigiai{
Số tròn chục liền trước của 320 là 310.
}
\end{ex}

\begin{ex}
Trong các số $485, 548, 458, 584$, số lớn nhất là số nào?
\shortans{584}
\loigiai{
So sánh hàng trăm rồi đến hàng chục: 584 là số lớn nhất.
}
\end{ex}

\begin{ex}
Trong các số $612, 621, 609, 619$, số bé nhất là số nào?
\shortans{609}
\loigiai{
Số bé nhất trong các số đã cho là 609.
}
\end{ex}

\begin{ex}
Cho các chữ số 2, 5, 9. Số lớn nhất có ba chữ số ghép từ cả ba chữ số này là số nào?
\shortans{952}
\loigiai{
Sắp xếp theo thứ tự giảm dần: 952.
}
\end{ex}

\begin{ex}
Cho các chữ số 0, 4, 7. Số bé nhất có ba chữ số khác nhau ghép từ ba chữ số này là số nào?
\shortans{407}
\loigiai{
Chữ số hàng trăm khác 0 nhỏ nhất là 4, tiếp theo là 0 và 7: 407.
}
\end{ex}

\begin{ex}
Tìm $x$ biết: $x$ là số liền trước của 780.
\shortans{779}
\loigiai{
Số liền trước của 780 là $780 - 1 = 779$.
}
\end{ex}

\begin{ex}
Tìm $y$ biết $y$ là số tròn chục nằm giữa 855 và 865.
\shortans{860}
\loigiai{
Số tròn chục duy nhất nằm giữa 855 và 865 là 860.
}
\end{ex}

\begin{ex}
Có bao nhiêu số tròn chục có hai chữ số?
\shortans{9}
\loigiai{
Các số đó là 10, 20, 30, 40, 50, 60, 70, 80, 90: có 9 số.
}
\end{ex}

\begin{ex}
Có bao nhiêu số tròn trăm từ 100 đến 1000?
\shortans{10}
\loigiai{
Các số: 100, 200, 300, 400, 500, 600, 700, 800, 900, 1000 (10 số).
}
\end{ex}

\begin{ex}
Viết số thích hợp vào chỗ chấm: $600 + 40 + 5 = \dots$
\shortans{645}
\loigiai{
$600 + 40 + 5 = 645$.
}
\end{ex}

\begin{ex}
Viết số thích hợp vào chỗ chấm: $900 + 8 = \dots$
\shortans{908}
\loigiai{
$900 + 8 = 908$.
}
\end{ex}

\begin{ex}
Số 354 gồm mấy chục nếu chỉ xét giá trị hàng chục?
\shortans{5}
\loigiai{
Chữ số hàng chục của 354 là 5.
}
\end{ex}

\begin{ex}
Điền số thích hợp: $299 < \dots < 301$.
\shortans{300}
\loigiai{
Số nằm giữa 299 và 301 là 300.
}
\end{ex}

\begin{ex}
Tổng của số lớn nhất có ba chữ số và số bé nhất có một chữ số là bao nhiêu?
\shortans{999}
\loigiai{
Số lớn nhất có ba chữ số là 999, số bé nhất có một chữ số là 0: $999 + 0 = 999$.
}
\end{ex}

\begin{ex}
Hiệu giữa số bé nhất có ba chữ số và số lớn nhất có hai chữ số là bao nhiêu?
\shortans{1}
\loigiai{
Số bé nhất có ba chữ số là 100, số lớn nhất có hai chữ số là 99: $100 - 99 = 1$.
}
\end{ex}

\begin{ex}
Từ ba chữ số 1, 2, 3 có thể lập được bao nhiêu số có ba chữ số khác nhau?
\shortans{6}
\loigiai{
Các số: 123, 132, 213, 231, 312, 321 (tổng cộng 6 số).
}
\end{ex}

\begin{ex}
Một đoàn tàu có 8 toa chở khách, mỗi toa có 100 chỗ ngồi. Hỏi đoàn tàu chở được tối đa bao nhiêu hành khách?
\shortans{800}
\loigiai{
$8 \times 100 = 800$ chỗ ngồi.
}
\end{ex}

\begin{ex}
Số nào cộng với 350 thì được kết quả là 500?
\shortans{150}
\loigiai{
$500 - 350 = 150$.
}
\end{ex}
